\documentclass[12pt,a4paper]{report}
\usepackage[utf8]{inputenc}
\usepackage{amsmath}
\usepackage{amsfonts}
\usepackage{amssymb}
\usepackage{amsthm}
\usepackage{hyperref}

\usepackage{multicol}
\usepackage{fancyhdr}
\usepackage[inline]{enumitem}
\usepackage{tikz}
\usepackage{tikz-cd}
\usetikzlibrary{calc}
\usetikzlibrary{shapes.geometric}
\usetikzlibrary{positioning}
\usepackage[margin=0.5in]{geometry}
\usepackage{xcolor}

\hypersetup{
    colorlinks=true,
    linkcolor=blue,
    filecolor=magenta,      
    urlcolor=cyan,
    pdftitle={Tensors},
    pdfpagemode=FullScreen,
    }

%\urlstyle{same}

\newcommand{\CLASSNAME}{Math 5050 -- Special Topics: Differential Geometry}
\newcommand{\STUDENTNAME}{Paul Carmody}
\newcommand{\ASSIGNMENT}{Section 23: Integration Manifolds }
\newcommand{\DUEDATE}{November 19, 2025}
\newcommand{\PROFESSOR}{Professor Berchenko-Kogan}
\newcommand{\SEMESTER}{Fall 2025}
\newcommand{\SCHEDULE}{TBD}
\newcommand{\ROOM}{Remote}

\newcommand{\MMN}{M_{m\times n}}
\newcommand{\FF}{\mathcal{F}}

\pagestyle{fancy}
\fancyhf{}
\chead{ \fancyplain{}{\CLASSNAME} }
%\chead{ \fancyplain{}{\STUDENTNAME} }
\rhead{\thepage}
\input{../MyMacros}

\newcommand{\RED}[1]{\textcolor{red}{#1}}
\newcommand{\BLUE}[1]{\textcolor{blue}{#1}}

\newcommand{\SUPP}{\operatorname{supp}}
\newcommand{\CLOSURE}{\operatorname{closure of}}
\newcommand{\BD}{\operatorname{bd}}
\newcommand{\HHN}{\mathcal{H}^n}

\begin{document}

\begin{center}
	\Large{\CLASSNAME -- \SEMESTER} \\
	\large{ w/\PROFESSOR}
\end{center}
\begin{center}
	\STUDENTNAME \\
	\ASSIGNMENT -- \DUEDATE\\
\end{center} 

\textbf{Key Points}

\HLINE
\begin{enumerate}[label=23.\arabic*]

\item \textbf{Area of an ellipse}

Use the change-of-variables formula to compute the area enclosed by the ellipse
\begin{align*}
	x^2/a^2+y^2/b^2=1
\end{align*}in $\R^2$.

\BLUE{Let $G(r,\theta) = (ar\cos\theta, br\sin\theta)$ be the transformation.  Then  the Jacobian of $G$
\begin{align*}
	dG &= \TWOXTWO{a\cos\theta}{-ar\cos\theta}{b\sin\theta}{br\sin\theta}.\\
	\det dG &= ab(r\cos^2\theta+r\sin^2\theta) = abr.
\end{align*}The change of variable formula states that
\begin{align*}
	A &= \int_0^{2\pi}\int_0^1 abr\,dr\, d\theta = ab(2\pi)\frac{1}{2} = \pi ab
\end{align*}
}

\item \textbf{Characterization of boundedness in $\R^n$}

Prove that a subset $A \subset \R^n$ is bounded if and only if its closure $\CONJ{A}$ in $\R^n$ is compact.

\BLUE{Since $A$ is bounded. Let $M$ be the maximum distance between any two points in $A$
\begin{align*}
	M &= \max_{x,y} \{x,y \in A\}
\end{align*}by defintiion, $\CONJ{A}$ is closed.  Hence, $\CONJ{A}$ is closed and bounded, which is compact.
}

\item \textbf{Integral under a diffeomorphism}

Suppose $N$ and $M$ are connected, oriented, $n$-manifolds and $F: N \to M$ is a diffeomorphism.  Prove that for any $\omega \in \Omega_c^k(M)$,
\begin{align*}
	\int_NF^*\omega = \pm \int_M \omega
\end{align*} where the sign depends on whether $F$ is orientation-preserving or orientation-reversing.

\item \textbf{Stoke's theorem}

Prove Stokes's theorem $\R^n$ and for $\mathcal{H}^n$.

\item \textbf{Area form on the sphere $S^2$}

Prove that the area form $\omega$ on $S^2$ in Example 23.11 is equal to the orientation form
\begin{align*}
	x\,dy\wedge dz-y \,dx\wedge dz+z\,dx\wedge dy
\end{align*}of $S^2$ in Problem 22.9.


\end{enumerate}
\end{document}
